
\section{在内部插入标记点}
\label{sec:3-8}

本节考察带若干内部标记点的伪全纯圆盘模空间，并据此构造算子
$\mathfrak p,\mathfrak q,\mathfrak r$。\liftIndex{I136}
这些算子可用于研究障碍类或弗洛尔上同调同限制同态
\[
 i^*\colon H^k(M)\longrightarrow H^k(L)
\]
以及 Gysin 同态
\[
 i_!\colon H^k(L)\longrightarrow H^{n+k}(M)
\]
之间的关系。本节结果还将用于研究第 6 章所构造谱序列的微分，并用于
定理 H、I 的证明。

\subsection{算子 $\mathfrak p$}
\label{subsec:3-8-1}

先回顾 Gysin 同态的定义：
\[
 i_!
 =\operatorname{PD}_M\circ i_*\circ\operatorname{PD}_L
 \colon H^k(L)\longrightarrow H^{n+k}(M),
\]
其中 $\operatorname{PD}$ 是 $M$ 与 $L$ 上的庞加莱对偶映射，
$i\colon L\to M$ 是自然嵌入。\liftIndex{I135}

令
\[
 \bigl(C(L;\Lambda_{0,\mathrm{nov}}),\mathfrak m\bigr)
\]
为定理 3.5.11 所构造的滤过 $A_\infty$ 代数，并令
$(S^*(M;\mathbb Z),\delta)$ 为 $M$ 的光滑奇异链复形。把余维为 $k$
的链视为次数为 $k$ 的上链后，我们把 $S^*(M;\mathbb Z)$ 看作上链复形。
对任意交换环 $R$，置
\[
 S^*(M;R)=S^*(M;\mathbb Z)\otimes R.
\]
Gysin 同态由下列上链映射诱导：
\[
 \begin{aligned}
 i_!\colon S^k(L;R)&\longrightarrow S^{k+n}(M;R),\\
 (|P|,f)&\longmapsto (|P|,i\circ f).
 \end{aligned}
\]

为了精确陈述主要技术结果——定理 3.8.9——还需要回顾一点
格罗莫夫--威腾理论。令
$0\ne\widetilde\beta\in\pi_2(M)$。考虑模空间
$\widetilde{\mathcal M}^{\mathrm{reg}}_{0,2}
(M;\widetilde\beta)$；它由所有
\[
 \bigl((S^2,(z_0^+,z_1^+)),w\bigr)
\]
组成，其中 $w\colon S^2\to M$ 是同伦类为 $\widetilde\beta$ 的伪全纯
映射，而 $(z_0^+,z_1^+)$ 是 $S^2$ 上一对互异点。用
$\mathrm{PSL}(2;\mathbb C)$ 的作用除
$\widetilde{\mathcal M}^{\mathrm{reg}}_{0,2}
(M;\widetilde\beta)$，再用稳定映射紧化所得商空间，并把最终模空间记作
$\mathcal M_{0,2}(M;\widetilde\beta)$。
\translatornote{原书在取商句中写
$\mathcal M_{0,2}(M;\widetilde\beta)$，但此前唯一新定义的是带波浪号与
$\mathrm{reg}$ 上标的参数化正则模空间；而无波浪号的记号紧接着才赋给
取商并紧化后的模空间。译文据此恢复取商对象。}
记号 $\mathcal M_{0,2}(M;\widetilde\beta)$ 中的下标 $0$ 表示伪全纯映射
定义域的亏格。在标记点 $z_1^+,z_0^+$ 处有两个评价映射，合写为
\[
 ev=(ev_1,ev_0)\colon
 \mathcal M_{0,2}(M;\widetilde\beta)\longrightarrow M^2.
\]
考虑 $M$ 中的链
\begin{equation}
 ev_{0*}\left(
 \mathcal M_{0,2}(M;\widetilde\beta)_{ev_1}
 \mathbin{\times_M}L
 \right).
\tag{3.8.1}\label{eq:3-8-1}
\end{equation}
依 \textbf{[FuOn99II]} 的构造，它带有仓西结构。把它的虚基本链记作
\[
 GW_{0,1}(M;\widetilde\beta)(L)
 \in S^*(M;\mathbb Q).
\]
这里的次数 $*$（即余维）为
\[
 *=n-2c_1(M)[\widetilde\beta]+2.
\]
更精确地说，我们固定
$\mathcal M_{0,2}(M;\widetilde\beta)_{ev_1}\mathbin{\times_M}L$
上的一个三角剖分以及顶点的一个次序，从而把
\eqref{eq:3-8-1} 视为 $M$ 中的光滑奇异链；见定义 A1.28。

于是定义一个 $\Lambda_{0,\mathrm{nov}}$-系数上链
\begin{equation}
\begin{aligned}
 GW_{0,1}(M)(L)
 &=
 \sum_{\widetilde\beta}
 T^{\omega[\widetilde\beta]}
 e^{c_1(M)[\widetilde\beta]}
 GW_{0,1}(M;\widetilde\beta)(L)\\
 &\in S^{n+2}(M;\Lambda_{0,\mathrm{nov}}).
\end{aligned}
\tag{3.8.2}\label{eq:3-8-2}
\end{equation}
\liftIndex{I134}\liftIndex{I138}

下面引入循环 bar 复形
$B_k^{\mathrm{cyc}}C[1]$。\liftIndex{I132}\liftIndex{I133}
\liftIndex{I137}
设 $C$ 是滤过 $A_\infty$ 代数，定义自同构
\[
 \operatorname{cyc}\colon B_kC[1]\longrightarrow B_kC[1]
\]
如下：
\[
 \operatorname{cyc}(x_1\otimes\cdots\otimes x_k)
 =
 (-1)^{
  \deg' x_k\left(\sum_{i=1}^{k-1}\deg' x_i\right)}
 x_k\otimes x_1\otimes\cdots\otimes x_{k-1}.
\]
它在 $B_kC[1]$ 上诱导一个 $\mathbb Z_k$-作用。以
$B_k^{\mathrm{cyc}}C[1]$ 表示 $\operatorname{cyc}$ 的不变部分，并置
\[
 B^{\mathrm{cyc}}C[1]
 =
 \widehat{\bigoplus}_{k}B_k^{\mathrm{cyc}}C[1],
\]
即这些不变部分的完备直和。

\begin{remark}{3.8.3}\label{rem:3-8-3}
文献中有时出现的循环 bar 复形概念，例如 \textbf{[GeJo90]} 中的概念，
是关于对称群 $\mathfrak S_k$ 在 $B_kC[1]$ 上的作用作对称化。这里则只
关于循环群 $\mathbb Z_k$ 作对称化。我们以 $E_kC$ 表示关于
$\mathfrak S_k$-作用的对称化；见定义 3.8.30。
\end{remark}

\begin{statement}{引理--定义}{3.8.4}\label{lemdef:3-8-4}
设 $(C,\mathfrak m)$ 是 $A_\infty$ 代数，
$(BC[1],\widehat d)$ 是相伴的 bar 复形，则
$B^{\mathrm{cyc}}C[1]$ 在映射 $\widehat d$ 下不变。因此定义
\[
 \widehat d^{\mathrm{cyc}}
 :=
 \left.\widehat d\right|_{B^{\mathrm{cyc}}C[1]}.
\]
\end{statement}

\begin{proof}
令 $x_1\otimes\cdots\otimes x_k\in B_kC[1]$，并考虑元素
\begin{equation}
 [x_1\otimes\cdots\otimes x_k]
 =
 \sum_{i=1}^{k}
 \operatorname{cyc}^{i}(x_1\otimes\cdots\otimes x_k)
 \in B_k^{\mathrm{cyc}}C[1].
\tag{3.8.5}\label{eq:3-8-5}
\end{equation}
于是
\begin{equation}
\begin{aligned}
 &\widehat d([x_1\otimes\cdots\otimes x_k])\\
 &\quad=
 \sum_{i\leq j+1}\sum_{\ell=1}^{k}
 (-1)^{
  \deg' x_\ell+\cdots+\deg' x_{\ell+i-2}+\alpha_\ell}\\
 &\qquad\qquad
 x_\ell\otimes\cdots\otimes
 \mathfrak m_{j-i+1}
 (x_{\ell+i-1},\ldots,x_{\ell+j-1})
 \otimes\cdots\otimes x_{\ell-1},
\end{aligned}
\tag{3.8.6}\label{eq:3-8-6}
\end{equation}
其中
\[
 \alpha_\ell
 =
 \sum_{a<\ell\leq b}\deg' x_a\,\deg' x_b.
\]
这里采用约定 $x_{k+m}=x_m$。容易验证，
\eqref{eq:3-8-6} 的右端属于 $B^{\mathrm{cyc}}C[1]$。

若 $\mathbb Q\subseteq R$，则形如
$[x_1\otimes\cdots\otimes x_k]$ 的元素张成
$B_k^{\mathrm{cyc}}C[1]$，引理在这种情形下已得证。

对一般的 $R$，这类元素未必张成 $B_k^{\mathrm{cyc}}C[1]$，因而需要换一种
论证。只需考虑系数环为 $\mathbb Z$ 的情形。假设
\[
 x_1\otimes\cdots\otimes x_k
 =
 \underbrace{\mathbf x\otimes\cdots\otimes\mathbf x}_{m'\text{ 次}},
 \qquad
 \mathbf x\in B_mC[1],
 \qquad
 mm'=k.
\]
则
\[
 \operatorname{cyc}^{m}(x_1\otimes\cdots\otimes x_k)
 =
 (-1)^{(m'-1)\deg'\mathbf x}
 x_1\otimes\cdots\otimes x_k,
 \qquad
 \deg'\mathbf x=\sum_{i=1}^{m}\deg' x_i.
\]
考虑
\[
 [x_1\otimes\cdots\otimes x_k]^{(m)}
 =
 \sum_{i=1}^{m}
 \operatorname{cyc}^{i}(x_1\otimes\cdots\otimes x_k).
\]
那么
\begin{equation}
\begin{aligned}
 &\widehat d\bigl([x_1\otimes\cdots\otimes x_k]^{(m)}\bigr)\\
 &\quad=
 \sum_{i\leq j+1}\sum_{\ell=1}^{m}
 (-1)^{
  \deg' x_\ell+\cdots+\deg' x_{\ell+i-2}+\alpha_\ell}\\
 &\qquad\qquad
 x_\ell\otimes\cdots\otimes
 \mathfrak m_{j-i+1}
 (x_{\ell+i-1},\ldots,x_{\ell+j-1})
 \otimes\cdots\otimes x_{\ell-1}.
\end{aligned}
\tag{3.8.7}\label{eq:3-8-7}
\end{equation}
由于该序列以 $m$ 为周期，即对所有整数 $r$ 有
$x_{m+r}=x_r$，故 \eqref{eq:3-8-7} 属于
$B^{\mathrm{cyc}}C[1]$。形如
$[x_1\otimes\cdots\otimes x_k]^{(m)}$ 的元素生成
$B_k^{\mathrm{cyc}}C[1]$，引理得证。（尽管这套代数论证对挠系数也成立，
目前我们只能在有理系数上证明主要几何结果定理 3.8.9。）
\end{proof}

\begin{remark}{3.8.8}\label{rem:3-8-8}
定义 $\widehat d^{\mathrm{cyc}}$ 不需要任何额外的对称性。庞加莱对偶配对
下
\[
 \bigl\langle x_0,\mathfrak m_k(x_1,\ldots,x_k)\bigr\rangle
\]
的循环对称性问题（见猜想 3.6.48），与定义算子
$\widehat d^{\mathrm{cyc}}$ 是性质不同的问题。
\end{remark}

本节的一个主要结果如下。

\begin{theorem}{3.8.9}\label{thm:3-8-9}
对任意包含 $\mathbb Q$ 的环 $R$，存在一列次数为 $n+1$ 的算子
\[
 \mathfrak p_k\colon
 B_k^{\mathrm{cyc}}C(L;\Lambda_{0,\mathrm{nov}})[1]
 \longrightarrow
 S^*(M;\Lambda_{0,\mathrm{nov}})
 \qquad(k=0,1,2,\ldots),
\]
具有下列性质。令
\[
 \mathfrak p\colon
 B^{\mathrm{cyc}}C(L;\Lambda_{0,\mathrm{nov}})[1]
 \longrightarrow
 S^*(M;\Lambda_{0,\mathrm{nov}})
\]
为限制在
$B_k^{\mathrm{cyc}}C(L;\Lambda_{0,\mathrm{nov}})[1]$ 上等于
$\mathfrak p_k$ 的算子。以 $\delta_M$ 表示
$S^*(M;\Lambda_{0,\mathrm{nov}})$ 上的
$\Lambda_{0,\mathrm{nov}}$-系数上边界算子，也就是 $M$ 上经典上边界
算子的显然延拓。则
\begin{align}
 \mathfrak p_1
 &\equiv i_!\pmod{\Lambda^+_{0,\mathrm{nov}}},
 \tag{3.8.10.1}\label{eq:3-8-10-1}\\
 \mathfrak p\circ\widehat d^{\mathrm{cyc}}
 +\delta_M\circ\mathfrak p_k
 &=0
 \qquad(k>0),
 \tag{3.8.10.2}\label{eq:3-8-10-2}\\
 \mathfrak p_1\circ\mathfrak m_0(1)
 +\delta_M\circ\mathfrak p_0(1)
 +GW_{0,1}(M)(L)
 &=0.
 \tag{3.8.10.3}\label{eq:3-8-10-3}
\end{align}
还可把 $\mathfrak p_k$ 延拓为
\[
 \mathfrak p_k^+\colon
 B_k^{\mathrm{cyc}}C(L;\Lambda_{0,\mathrm{nov}})^+[1]
 \longrightarrow
 S^*(M;\Lambda_{0,\mathrm{nov}}),
\]
把同伦幺元纳入其中；以
$\mathfrak p^+,\widehat d^{+\mathrm{cyc}}$ 分别代替
$\mathfrak p,\widehat d^{\mathrm{cyc}}$ 后，
\eqref{eq:3-8-10-2} 仍然成立。这里
$\widehat d^{+\mathrm{cyc}}$ 由 $\mathfrak m^+$ 按照
$\widehat d^{\mathrm{cyc}}$ 的同样方式定义。此外，
\begin{align}
 \mathfrak p_k^+
 [\mathbf e^+,x_1,\ldots,x_{k-1}]
 &=0
 \qquad(k\ne1,2),
 \tag{3.8.10.4}\label{eq:3-8-10-4}\\
 \mathfrak p_1^+(\mathbf e^+)
 &=\operatorname{PD}([L]),
 \tag{3.8.10.5}\label{eq:3-8-10-5}\\
 \mathfrak p_2^+(\mathbf e^+,x)
 &=x.
 \tag{3.8.10.6}\label{eq:3-8-10-6}
\end{align}
\end{theorem}
\liftIndex{I136}

这里
\[
 C(L;\Lambda_{0,\mathrm{nov}})^+
 =
 C(L;\Lambda_{0,\mathrm{nov}})
 \oplus\Lambda_{0,\mathrm{nov}}\mathbf e^+
 \oplus\Lambda_{0,\mathrm{nov}}\mathbf f
\]
如第 3.3 节，而 $[L]$ 是 $L$ 的基本闭链。

定理 3.8.9 的证明将在第 3.8.3 小节与第 7.4 节给出；证明
\eqref{eq:3-8-10-4}--\eqref{eq:3-8-10-6} 还需要第 7.3 节的论证。
关于 \eqref{eq:3-8-10-3} 所涉定向问题，见第 8.10.2 小节的
命题 8.10.6。

\subsection{障碍类 $o_k(L)$ 的消失及其应用}
\label{subsec:3-8-2}

本小节讨论定理 3.8.9 的若干应用。令
\[
 [GW_{0,1}(M)(L)]
 \in H^*(M;\mathbb Q)\otimes\Lambda_{0,\mathrm{nov}}
\]
为 $GW_{0,1}(M)(L)$ 的上同调类。\liftIndex{I138}
在下一个定理中，$\operatorname{PD}_M$ 表示庞加莱对偶
\[
 \operatorname{PD}_M\colon
 H_*(M;\mathbb Q)\longrightarrow H^{2n-*}(M;\mathbb Q).
\]

\begin{theorem}{3.8.11}\label{thm:3-8-11}
设 $L\subset M$ 是相对 spin 拉格朗日子流形。

\smallskip
\noindent\textup{(3.8.12.1)}
若 $[GW_{0,1}(M)(L)]\ne0$，则
$\widehat{\mathcal M}(L)=\varnothing$。

\noindent\textup{(3.8.12.2)}
若 $[GW_{0,1}(M)(L)]=0$，则障碍类 $[o_i(L)]$ 属于 Gysin 同态
\[
 i_!\colon H^{2-\mu_i}(L;\mathbb Q)
 \longrightarrow H^{n+2-\mu_i}(M;\mathbb Q)
\]
的核。

\noindent\textup{(3.8.12.3)}
若 $[GW_{0,1}(M)(L)]$ 不具有
$c\operatorname{PD}_M[L]$ 的形式，则
$\widehat{\mathcal M}_{\mathrm{weak}}(L)=\varnothing$。

\noindent\textup{(3.8.12.4)}
若
$[GW_{0,1}(M)(L)]=c\operatorname{PD}_M[L]$，则
\translatornote{原书此项把 $o_i(L;\mathrm{weak})$ 直接写成 Gysin 同态之核
中的元素，漏去上同调类方括号；同一定理的 (3.8.12.2)、证明以及
推论 3.8.16 均明确使用相应的上同调类，译文据此补回。}
\[
 [o_i(L;\mathrm{weak})]
 \in
 \operatorname{Ker}\left(
 i_!\colon H^{2-\mu_i}(L;\mathbb Q)
 \longrightarrow H^{n+2-\mu_i}(M;\mathbb Q)
 \right).
\]

\noindent\textup{(3.8.12.5)}
设 $[GW_{0,1}(M)(L)]=c\operatorname{PD}_M[L]$，
$b\in\widehat{\mathcal M}_{\mathrm{weak}}(L)$，并且
$[L]\ne0$ 于 $H(M;\mathbb Q)$ 中，则
\[
 c=\mathfrak{PO}(b).
\]
\end{theorem}

\begin{proof}
回顾 $G(L)$ 是由
\[
 G_0(L)
 =
 \left\{
 (\omega[\beta],\mu(\beta))
 \ \middle|\
 \beta\in\Pi(L),\ 
 \mathcal M(L;\beta)\ne\varnothing
 \right\}
\]
生成的幺半群。取满足条件 3.1.6 的幺半群 $G\supset G(L)$，并写成
\[
 G=\{(\lambda_i,\mu_i)\mid i=0,1,2,\ldots\},
\]
其中
\[
 \lambda_i\leq\lambda_{i+1},
 \qquad
 \lambda_i=\lambda_{i+1}\Longrightarrow\mu_i<\mu_{i+1}.
\]
为简化记号，假设对每个 $i$ 都有
$\lambda_i<\lambda_{i+1}$；其他情形只需作显然的修改即可处理，见
第 3.6.2 小节定理 3.6.18 的证明。

假设 $j=1,\ldots,i-1$ 时 $[o_j(L)]=0$。这意味着可以选取
\[
 b_{(i-1)}
 =
 \sum_{j=1}^{i-1}
 T^{\lambda_j}e^{\mu_j/2}b_j,
\]
使
\[
 \mathfrak m(e^{b_{(i-1)}})
 \equiv0
 \pmod{T^{\lambda_{i-1}}\Lambda^+_{0,\mathrm{nov}}}.
\]
还要注意
\[
 \widehat d(e^{b_{(i-1)}})
 \equiv
 T^{\lambda_i}o_i(L)
 \pmod{T^{\lambda_i}\Lambda^+_{0,\mathrm{nov}}}.
\]
此外，$e^{b_{(i-1)}}$ 在 $\operatorname{cyc}$ 作用下不变，并且
\[
 \widehat d^{\mathrm{cyc}}(e^{b_{(i-1)}})
 =
 e^{b_{(i-1)}}\otimes
 \mathfrak m(e^{b_{(i-1)}})\otimes
 e^{b_{(i-1)}}.
\]

现在利用 \eqref{eq:3-8-10-2} 与 \eqref{eq:3-8-10-3} 计算，得到
\[
 \mathfrak p\left(
 \widehat d^{\mathrm{cyc}}(e^{b_{(i-1)}})
 \right)
 =
 -\delta_M\left(\mathfrak p(e^{b_{(i-1)}})\right)
 -GW_{0,1}(M)(L).
\]
再用 \eqref{eq:3-8-10-1}，有
\begin{equation}
\begin{aligned}
 T^{\lambda_i}i_!(o_i(L))
 &\equiv
 -\delta_M\left(\mathfrak p(e^{b_{(i-1)}})\right)
 -GW_{0,1}(M)(L)\\
 &\hspace{12em}
 \pmod{T^{\lambda_i}\Lambda^+_{0,\mathrm{nov}}}.
\end{aligned}
\tag{3.8.13.1}\label{eq:3-8-13-1}
\end{equation}
由于 $GW_{0,1}(M)(L)$ 给出 $M$ 上的一个
$\Lambda_{0,\mathrm{nov}}$-系数闭上链，
\eqref{eq:3-8-13-1} 在上同调类上给出
\begin{equation}
 T^{\lambda_i}i_!([o_i(L)])
 \equiv
 -[GW_{0,1}(M)(L)]
 \pmod{T^{\lambda_i}\Lambda^+_{0,\mathrm{nov}}}.
\tag{3.8.13.2}\label{eq:3-8-13-2}
\end{equation}
所以，若所有 $i$ 都有 $[o_i(L)]=0$，则对所有 $i$，
\[
 [GW_{0,1}(M)(L)]
 \equiv0
 \pmod{T^{\lambda_i}\Lambda^+_{0,\mathrm{nov}}},
\]
从而 $[GW_{0,1}(M)(L)]=0$。反过来，若
$[GW_{0,1}(M)(L)]=0$，则由 \eqref{eq:3-8-13-2} 得
$i_!([o_i(L)])=0$。

关于 $[o_i(L;\mathrm{weak})]$ 的结论可用
\eqref{eq:3-8-10-4} 与 \eqref{eq:3-8-10-5} 作类似证明。定理证毕。
\end{proof}

定理 3.8.11 有若干具体推论。先从下列引理开始。

\begin{lemma}{3.8.14}\label{lem:3-8-14}
若
\[
 [GW_{0,1}(M;\widetilde\beta)(L)]
 \in H^*(M;\mathbb Q),
 \qquad *=2n\ \text{或}\ 2n-1,
\]
则
$[GW_{0,1}(M;\widetilde\beta)(L)]=0$。
\end{lemma}

\begin{proof}
由
\[
 \dim\mathcal M_{0,2}(M;\widetilde\beta)
 =
 2n+2c_1(M)[\widetilde\beta]-2
\]
可得
\[
 *=n-2c_1(M)[\widetilde\beta]+2.
\]
与 $\mathcal M_{0,2}(M;\widetilde\beta)$ 类似定义的
$\mathcal M_{0,1}(M;\widetilde\beta)$，其虚维数因而为
\[
 2n-4+2c_1(M)[\widetilde\beta]
 =
 3n-*-2.
\]
所以由假设，
\[
 \dim\mathcal M_{0,1}(M;\widetilde\beta)
 =n-2\quad\text{或}\quad n-1.
\]
于是纤维积
\[
 \mathcal M_{0,1}(M;\widetilde\beta)_{ev_0}
 \mathbin{\times_M}L
\]
为空。可以选取与遗忘映射相容的
$\mathcal M_{0,2}(M;\widetilde\beta)$ 上的扰动，从而得到遗忘映射
\[
 \mathcal M_{0,2}(M;\widetilde\beta)
 \longrightarrow
 \mathcal M_{0,1}(M;\widetilde\beta),
\]
并使 $ev_0$ 与它相容。
\translatornote{原书先正确称此映射为 forgetful map，紧接着却把同一个
模空间间的映射称作 evaluation map；其定义是忘掉一个标记点，译文统一为
“遗忘映射”。}
见 \textbf{[FuOn99II]} 第 23 节。因此可得
$[GW_{0,1}(M;\widetilde\beta)(L)]=0$。
\end{proof}

\begin{remark}{3.8.15}\label{rem:3-8-15}
这里的论证比第 7.3 节中构造同伦幺元时的类似论证简单。这是因为处理
格罗莫夫--威腾不变量时可以在上同调层次工作，不必在链层次工作。
\end{remark}

\begin{corollary}{3.8.16}\label{cor:3-8-16}
若 $[o_i(L)]$ 已定义且属于 $H^n(L;\mathbb Q)$，则
$[o_i(L)]=0$。若 $[o_i(L)]$ 属于 $H^{n-1}(L)$，则
$i_!([o_i(L)])=0$。同样的结论也适用于
$[o_i(L;\mathrm{weak})]$。
\end{corollary}

\begin{proof}
注意
\[
 i_!\colon H^n(L)\longrightarrow H^{2n}(M)
\]
是单射。结论由定理 3.8.11 与引理 3.8.14 立即得到。
\end{proof}

\begin{corollary}{3.8.17}\label{cor:3-8-17}
设 $\dim_{\mathbb R}M=6$、$c_1(M)=0$，并设 $L$ 是相对 spin
拉格朗日子流形。若
\[
 i_*\colon H_*(L;\mathbb Q)\longrightarrow H_*(M;\mathbb Q)
\]
是单射，则 $L$ 无障碍。
\end{corollary}

这由定理 3.8.11 与引理 3.8.14 立即得到。下一个推论也可由
推论 3.8.16 与定理 3.6.43 得到。

\begin{corollary}{3.8.18}\label{cor:3-8-18}
任何作为有理同调球面的相对 spin 拉格朗日子流形都是弱无障碍的。
\end{corollary}

还要指出，定理 3.8.11 也给出了格罗莫夫--威腾不变量与 $L$ 的同调类
之间关系的一些新信息。例如有：

\begin{corollary}{3.8.19}\label{cor:3-8-19}
设 $M$ 是满足 $c_1(M)=0$ 的辛流形，$L$ 是马斯洛夫类为零的相对 spin
拉格朗日子流形。若 $H^2(L;\mathbb Q)=0$，则
\[
 [GW_{0,1}(M)(L)]=0.
\]
\end{corollary}

这显然由定理 3.8.11 得到。例如，当 $n=4$ 时，这意味着
$H_4(M;\mathbb Z)$ 中某些中维同调类不能由八维辛流形 $M$ 内、满足
$c_1(M)=0$ 的相对 spin 拉格朗日有理同调 $4$-球面实现。要得到更明确的
结果，还需要计算格罗莫夫--威腾不变量。

\begin{remark}{3.8.20}\label{rem:3-8-20}
在本书 2000 年的预印本版本 \textbf{[FOOO00]} 中，我们曾断言：若
\[
 H_*(L;\mathbb Q)\longrightarrow H_*(M;\mathbb Q)
\]
是单射，则 $L$ 无障碍。这个断言似乎并不正确；我们漏掉了
\eqref{eq:3-8-10-3} 右端出现的项 $GW_{0,1}(M)(L)$。下一个小节将解释
为何会出现这一项。不过，利用将在第 3.8.4、3.8.9 小节引入的算子
$\mathfrak q$ 与 $\mathfrak p_{\ell,k}$，可以部分克服这一困难；例如见
推论 3.8.43。特别地，\textbf{[FOOO00]} 中的“定理 B”（即本书的
定理 I）仍可按 \textbf{[FOOO00]} 中原来的陈述得到证明；见第 6.5.2
小节。
\end{remark}

第 6 章还将利用算子 $\mathfrak p$ 证明：定理 D 中谱序列微分的像属于
Gysin 同态
\[
 i_!\colon H^*(L;\mathbb Q)\longrightarrow H^{*+n}(M;\mathbb Q)
\]
的核。
